\section{Why CoT Helps: A Time-Indexed Causal Chain}
\label{sec:synthesis}

\subsection{The proposed explanatory chain}

The central result cannot be ``an ablation lowers accuracy.'' The proposed
mechanism is a sequence of registered mediators:
\begin{equation}
  I_{{\rm route},k}
  \longrightarrow
  \pi_k
  \longrightarrow
  \Gamma_k
  \longrightarrow
  (\pi_{k+1},\textsc{stop})
  \longrightarrow
  Z
  \longrightarrow
  Y_{\rm full},
\label{eq:causal-chain}
\end{equation}
with side interventions
\[
  I_{{\rm state},k}\longrightarrow\Gamma_k,
  \qquad
  I_{\rm terminal}\longrightarrow Z.
\]
\(\pi_k\) is the aligned source occurrence actually retrieved;
\(\Gamma_k\) contains semantic progress and any visible/KV ledger;
\(Z\) is the registered terminal state; and \(Y_{\rm full}\) is the complete
registered continuation, parsed count, exactness, and absolute error.

Each intervention has a proximal claim gate. Corrupting a targeted-retrieval
component at \(Q_k\) must first change the probability or identity of the gold
next source. Patching \(P_k\) must first change the next source, repetition, or
stopping. Patching \(Z\) may change the full answer without changing coverage
and therefore isolates terminal readout. Clean-output rescue is evaluated
beside corruption. Associations among attention, state quality, and accuracy
are descriptive; causal language is reserved for paired activation
interventions and their randomized controls.

The contrast with Direct is computational. Direct may use the shared
predicate-conditioned record states through \(H_{\rm run}\), or it may make
many record contributions influence a late state through \(H_{\rm agg}\).
CoT exposes a recurrent schedule in which each source decision has a local,
testable endpoint and its result persists in the generated prefix. The claim
is an algorithmic path change only if the corresponding route, progress, and
terminal links in \cref{eq:causal-chain} are all supported.

\subsection{Noise has a measurable referent}

We do not use ``noise'' as a synonym for diffuse attention or an unattractive
PCA plot. In the main text it refers to one of two quantities:
\begin{enumerate}
  \item functional retrieval/update error---missed valid occurrences, false
  positives, repeated retrieval of one occurrence, wrong next source, or
  premature/late stopping; or
  \item nuisance-conditioned state dispersion
  \(\nu_{\rm state}\) in \cref{eq:state-snr}.
\end{enumerate}
Attention entropy, target mass, and selectivity are routing descriptors.
Parse, format, and truncation are behavioral failure categories. They are not
folded into one latent ``noise'' variable.

CoT can improve counting by reducing functional retrieval error: targeted
queries increase the margin for one unvisited source, while the trace makes
visited identities available for duplicate avoidance and stopping. This does
not predict immunity to scale. Larger \(N\) creates more opportunities for
misses, repeats, and stop errors, and longer \(L\) may reduce source margins.
The shape of this accumulation is estimated from the error decomposition in
\cref{eq:error-budget}; it is not imposed as an independent-step product law.

\subsection{Factorial intervention on route, state, and terminal readout}

On locked Qwen families, we cross clean/corrupt/rescued levels for the route,
progress carrier, and terminal state:
\begin{itemize}
  \item \textbf{Route intervention.} Ablate or wrong-source patch the frozen
  targeted component at \(Q_k\). First measure gold/emitted identity and
  unique coverage.
  \item \textbf{State intervention.} Patch a single clean \(P_k\) donor into a
  matched corrupted prefix, or separately transport trace/KV state. First
  measure next source, repeat, continue, and stop.
  \item \textbf{Terminal intervention.} Patch a single clean \(Z\) donor at
  the same \texttt{Total:} landmark. First measure complete numeric-string
  likelihood and then free-running exact count/MAE.
\end{itemize}
The decisive pattern is ordered degradation and rescue: route corruption
changes source identity/coverage; state corruption changes future retrieval
or stopping; the resulting terminal state and full answer change; restoring a
downstream clean mediator rescues only its descendants. A terminal patch that
restores the answer without restoring coverage is expected and identifies a
readout effect rather than a retrieval repair.

For an endpoint \(m\), we report paired raw differences and compute
\[
  \operatorname{Rec}(m)=
  \frac{m_{\rm rescue}-m_{\rm corrupt}}
       {m_{\rm clean}-m_{\rm corrupt}}
\]
only above a validation-frozen denominator gate. Intervals are clustered by
stimulus family and paired randomization tests respect the sibling structure.
No single normalized recovery ratio replaces the raw semantic endpoints.

\begin{figure}[H]
\centering
\fbox{\parbox[c][5.6cm][c]{0.94\textwidth}{
\centering
\textbf{TODO Figure: Integrated causal chain}\\[0.5em]
\(\boxed{\text{route}}\rightarrow
\boxed{\text{source identity / coverage}}\rightarrow
\boxed{\text{progress / next or stop}}\rightarrow
\boxed{\text{terminal state}}\rightarrow
\boxed{\text{complete answer}}\)\\[0.8em]
\resizebox{0.91\linewidth}{!}{\begin{tabular}{lccccc}
Condition & Gold identity & Unique coverage & State quality & Terminal & Exact count\\
\hline
Clean & high & high & high & correct & high\\
Route corrupt & \(\downarrow\) & \(\downarrow\) & \(\downarrow\) & damaged & \(\downarrow\)\\
+ clean progress patch & unchanged/partial & partial rescue & rescued & partial & partial\\
+ clean terminal patch & unchanged & unchanged & unchanged & rescued & rescued
\end{tabular}}\\[0.7em]
\small Show proximal endpoints before downstream answers, raw paired effects,
family-cluster intervals, and every matched control for Qwen3-8B and
Gemma4-E4B.
}}
\caption{Planned decisive result figure. A causal explanation requires
intervention effects to propagate through the registered semantic chain, not
merely an accuracy drop.}
\label{fig:mediation-todo}
\end{figure}

\begin{todobox}
\textbf{M1---integrated Qwen causal chain.}
On locked Qwen3-8B families, run the route-by-state-by-terminal factorial with
clean, corrupt, rescue, and sham levels. For every intervention, report the
proximal endpoint first: gold/emitted source identity for route; next
source/repeat/stop for progress; complete answer for terminal. Then report
unique coverage, false positives, duplicates, sequence likelihood, terminal
state quality, exactness, and MAE with family-cluster intervals and paired
randomization tests. Include single-donor natural patches, pattern/value/head
and residual path decomposition, random/same-layer/wrong-row/wrong-source
controls, and free-running validation before and after the first trace error.
Do not replace the chain by one accuracy drop or a cross-sectional mediation
regression.
\end{todobox}

\subsection{Common-evidence readout and trace-content controls}

Retrieval and readout must be separated without pretending that a
teacher-forced list is a natural Direct computation. We give both
reasoning-policy conditions the same correct, passage-ordered evidence ledger
and compare the state at the identical \texttt{Total:} token. This
\emph{common-evidence readout} intervention fixes unique coverage and is an
auxiliary control: if the interface gap disappears, the natural advantage is
consistent with retrieval mediation; if it remains, state construction or
terminal readout contributes after evidence is fixed.

An equal-length null scratchpad controls for extra tokens and recurrent
computation. Counterfactual traces preserve token length and grammar while
inserting a matched non-target, a repeated occurrence, or a false stop. A
trace-length-matched correct prefix controls for merely having more KV
positions. We also compare the same faithful ledger in a user prefix versus an
assistant prefix, shuffled faithful evidence, and Direct with a prefilled
ledger. These interventions test whether semantic trace content, rather than
verbosity, label tokens, or additional compute, drives the benefit.

\begin{todobox}
\textbf{M2---controls and Gemma functional replication.}
Repeat every decisive link in Gemma4-E4B using independently discovered but
identically defined functional components. Run the common-evidence readout,
equal-length null scratchpad, matched non-target, repeated-source, false-stop,
shuffled-ledger, and user-versus-assistant-prefix controls. Claim functional
replication if the same semantic links hold even when layers and head
identities differ. If only Qwen supports the complete chain, state the
model-specific boundary. Claim an improvement beyond coverage only if a
fixed-source-set intervention still changes the registered terminal state and
the complete answer.
\end{todobox}

\subsection{Claim ladder and falsifiers}

\begin{table}[H]
\centering
\small
\caption{Evidence required for each level of the paper's claim. Passing an
earlier level does not license a later conclusion.}
\label{tab:claim-ladder}
\begin{tabularx}{\textwidth}{@{}P{.18\textwidth}P{.31\textwidth}Y@{}}
\toprule
Claim level & Required evidence & Principal falsifier or weaker conclusion \\
\midrule
Behavior
  & trace-generating interfaces outperform Direct under matched stimuli
  & effect disappears under provenance, policy, or budget control \\
Process
  & source retrieval, duplicates, stopping, and Direct total formation are
  separately measured
  & only the final number is available \\
Representation
  & grouped held-out decoding and the registered dynamic profile
  & only token, index, position, or density decodes \\
Component causality
  & local semantic damage and clean single-donor rescue on frozen candidates
  & heatmap/ablation changes without the proximal endpoint \\
Integrated causality
  & effects propagate through coverage, progress, terminal state, and answer
  & null compute matches the gain or downstream restoration fails \\
Scope
  & functional chain replicates in both mechanism models and a compatible
  synthetic setting
  & implementation or result is checkpoint/template specific \\
\bottomrule
\end{tabularx}
\end{table}
